\section{Design Implications}
\subsection{Preserve a reliable backbone, then test a specific change}
The four-object exchange study gives a constructive answer to the paper's question. Receiver ordering provides a useful starting space; a few conditionally evaluated replacements substantially improve the complete material step without increasing its current dimension. The 35 resolved improvements and the median 65.2\% object-level gap reduction demonstrate local design value across both tested material representations. The operative quantity is the gain of the displacement produced by the modified model. Neither its direct radiation nor its transfer norm supplies that sign.

This design preserves what the receiver model already does well. A candidate is evaluated relative to the remaining currents, including their feedback and mutual coupling. After accepting a move, the current endpoint step, cached output, and conditional factors are updated before choosing another. The method consequently uses physical structure to decide a bounded model modification, rather than treating a current ranking as independent of the selected space or the present residual.

\subsection{Candidate coverage is an actionable bottleneck}
After summing one-decision gains over each object's nine common receiver checkpoints, the richer-reference score captures 98.1--99.6\% of the best gain in the declared neighborhood. Yet the twelve-incoming pool contains only 35.3\% of that opportunity for the near-contact object and 65.9\% for the asymmetric object. Improving proposal coverage is therefore a more direct next design task than replacing an already accurate full-dictionary score with a more elaborate regressor. The six learned-priority runs reinforce this reading: all trail the analytical anchor top-two rule on the heldout object.

Coverage and conditional interaction must be considered together. The positive-pair Maxwell witnesses show that two unfavorable singleton additions can jointly help. These witnesses motivate limited block proposals when a singleton screen is inconclusive; they do not supply a universal greedy or submodularity guarantee. The complete teacher comparison remains local to its declared dictionary and action horizon.

\subsection{Information structure and update utility are distinct}
Increasing a model's information volume can enlarge its curvature while changing its residual pullback in the wrong direction. The resolved replacement in Fig.~\ref{fig:information-pair} increases information volume and effective dimension and slightly improves the measured curvature fidelity, yet worsens the complete material step. This is an electromagnetic example of why both terms in Eq.~\eqref{eq:information-offset} are needed. Information diagnostics explain which directions change; signed full-quadratic gain decides whether the current task benefits.

All information comparisons use fixed physical material coordinates and the declared scale. Statistical Fisher or posterior interpretations additionally require an appropriate noise and prior model. The present normalized data scale is a numerical convention. The Gaussian full-reference spectrum has no weak direction at the specified threshold, while the voxel projection samples only sixteen material directions. These findings guide the weak-material gate without establishing a global voxel posterior model.

\subsection{Accuracy, certification, and cost have separate contracts}
Directional checks make a candidate decision inexpensive relative to constructing a full Jacobian: once $Jx$ is cached, a finalist needs $Jd$, without a new complete adjoint. The measured warm conditional selections take 0.594--3.400~s in the two representative physical timing states. They cost more than the cached one-pass control, and they buy a different level of full-objective checking and conditional search. Independent cold-deployment timing and net reconstruction savings remain unmeasured.

High-fidelity checks are also distinct from deterministic output bounds. Every accepted directed-verified exchange passes the frozen complete-gain numerical rule, but the implementation lacks a certified output-error radius for all candidates. Its terminal states therefore remain unresolved or move-budget stops. A rigorous neighborhood stopping claim requires upper bounds for every feasible move in that neighborhood, including those excluded by a shortlist.

The earlier structural-repair controls sharpen this distinction. Paired correction beats the specified direct enlargement, while equal-rank material Krylov correction is stronger in all recorded comparisons. The paid paired warm route costs 1.416 times its PCG control. Already-paid current histories can reduce repeated-solve marginal cost, but dedicated history acquisition removes the measured savings within eight uses. These results favor spending current-side information where it improves a declared model decision, with a standard material-space control and a complete acquisition ledger.

\subsection{From local design to an imaging platform}
The feedback and material-pullback identities apply to near- or far-field reception through the corresponding observation map. Known invertible calibration, consistently applied to the data, Jacobian, and noise covariance, preserves the whitened normal equation. For a raw Jacobian $F$ and calibration $T$,
\[
(TF)^*(T\Sigma_nT^*)^{-1}(TF)=F^*\Sigma_n^{-1}F.
\]
Finite aperture and discarded polarization instead remove channels. Amplitude correction cannot generally recover their material information. These synthetic design conditions connect the mechanism to full-wave extensions of far-field imaging platforms; measured hardware calibration is a separate validation task.

Accepted nonlinear updates also require constraints and a line search. The unit-step truth diagnostic shows why: 19 of 36 steps reduce error relative to the current material, 17 increase it, and 28 violate material constraints. It tests a different quantity from complete-quadratic fidelity. The successful frozen exchange results provide a specific basis for that next validation, rather than an already established reconstruction guarantee.

\section{Conclusion}
A current mode matters when material excitation can enter it and its interacting response can return to the receivers. Retained-space elimination exposes that response, and paired material pullbacks describe the complete change in a regularized material step. The resulting curvature change is not generally monotone because a current-space modification changes the model of existing measurements.

Receiver-centered conditional exchange turns this mechanism into a practical frozen-update design: at fixed current dimension, a few checked replacements improve 35 of 36 vector-Maxwell cases, with a median object-summed gap reduction of 65.2\%. The strongest diagnostic result is equally useful: accurate full-dictionary scoring coexists with incomplete shortpool coverage. Future improvement should target the physical proposal domain and its acquisition cost. Information-volume and pair-interaction witnesses specify why norm-only rankings are insufficient; directed gain checks specify what a useful replacement must actually accomplish.
